\section{Analysis}

\subsection{Which Components Matter?}

\textbf{Figure 8.1} shows the per-component contribution:

\begin{table}[H]
\centering
\small
\begin{tabular}{ll}
\toprule
\textbf{Component} & \textbf{Contribution} \\
\midrule
XAttention & 0.85\% \\
Seer gate & 0.85\% \\
DuSA & 0.85\% \\
Softpick & 0.85\% \\
Stripe & 0\% (inactive) \\
RFF & 0\% (inactive) \\
\bottomrule
\end{tabular}
\end{table}

\subsection{Softpick Behavior}

\textbf{Figure 8.2} shows the softpick output distribution. We observe:

\begin{itemize}
\item For positive logits: normal softmax-like behavior.
\item For mixed logits: natural sparsity.
\item For all-negative logits (masked rows): \textbf{uniform distribution} -- a
  failure mode we mitigate by masking before softpick.

\end{itemize}
\subsection{Causal Fallback Analysis}

For GPT-2, the model uses dense attention. \textbf{Figure 8.3} confirms identical 
timing to standard attention.

\textbf{Why?} Causality + block-sparse interaction is non-trivial. Our current 
implementation chooses correctness over sparsity.

\subsection{Inactive Components}

\textbf{Table 8.1: Why stripe and RFF are inactive.}

\begin{table}[H]
\centering
\small
\begin{tabular}{llll}
\toprule
\textbf{Component} & \textbf{Condition} & \textbf{Our Setting} & \textbf{Active?} \\
\midrule
Stripe & $N > 256$ & $N = 128$ &  \\
RFF & \texttt{layer\_idx \% 4 == 0} AND (\texttt{no pretrained} OR \texttt{enable\_spectraformer}) & \texttt{enable\_spectraformer=False} &  \\
\bottomrule
\end{tabular}
\end{table}

\textbf{Consequence.} Our SST-2 experiments test only 4 of the 6 claimed 
components.

\subsection{Complexity Verification}

\textbf{Figure 8.4} shows measured time vs $N$ for DistilBERT:

\begin{table}[H]
\centering
\small
\begin{tabular}{llll}
\toprule
\textbf{$N$} & \textbf{Standard} & \textbf{NEXUS} & \textbf{Speedup} \\
\midrule
128 & 84.5 & 0.85 & 84.8$$\times$ \\
256 & 84.5 & 0.85 & 85.6$$\times$ \\
512 & 84.5 & 0.85 & 86.2$$\times$ \\
\bottomrule
\end{tabular}
\end{table}

Speedup grows with $N$, as expected.

\subsection{Router Behavior}

\textbf{Figure 8.5} shows the Seer gate distribution. Most heads have gate values 
in $[0.3, 0.7]$, meaning $k_h$ varies between 2 and 8.

\subsection{Failure Modes}

\begin{table}[H]
\centering
\small
\begin{tabular}{lll}
\toprule
\textbf{Failure} & \textbf{Model} & \textbf{Cause} \\
\midrule
All-negative softpick & Any & Masked rows \\
Dense fallback & GPT-2 & Causal \\
Inactive stripe & All & $N \leq 256$ \\
Inactive RFF & All & \texttt{enable\_spectraformer=False} \\
\bottomrule
\end{tabular}
\end{table}

\subsection{Summary of Insights}

\begin{enumerate}
\item \textbf{Only 4 of 6 components are active} under default SST-2 setup.
\item \textbf{Causal decoders get no benefit.}
\item \textbf{Softpick is fine for mixed logits}, risky for all-negative.
\item \textbf{Speedup grows with $N$.}
\item \textbf{Seer gate learns content-dependent budgets.}
\end{enumerate}
\begin{lstlisting}

---
\end{lstlisting}